\chapter{The \texttt{phonon.solver} Package}
\label{sec:phonon_solver_pkg}

The \texttt{phonon.solver} package is the dense reference implementation
of the non-equilibrium Green's function (NEGF) plus self-consistent
Born approximation (SCBA) machinery used for the anharmonic phonon
transport calculations of this report.  It serves three roles
simultaneously:

\begin{enumerate}
  \item It is the canonical reference against which the production
        block-sparse / MPI solver in
        \texttt{quatrex.phonon.sse\_phonon\_phonon} is benchmarked.
        Both solvers share the same bubble kernel
        (\texttt{bubble.py}); cross-agreement is locked by a unit
        test that asserts agreement to $\sim\!10^{-12}$ (absolute
        tolerance) on a small synthetic Green's function.
  \item It is the driver behind the cutoff-sensitivity and
        FC3-compression-vs-transport sweeps that produce the
        results-section figures.  The cell that exposes those
        analyses lives in \texttt{finite\_analysis.transport\_quality}
        and feeds the solver via
        \texttt{phonon.solver.transmission\_q(...,
        M\_stacked\_override=...)}.
  \item It exposes the SCBA fixed-point loop
        (\texttt{scba\_loop}) and the bare bubble kernel
        (\texttt{bubble\_dense}) as library entry points reusable by
        any caller that wants to plug in a custom self-energy.
\end{enumerate}

This appendix walks through the package module by module.  Public
entry points are listed in \texttt{phonon/solver/\_\_init\_\_.py}; the
internal modules are described in the order in which they are pulled
in by the top-level drivers.

\section{Unit system and frequency grid}
\label{sec:solver_units_grid}

The solver operates internally in \textbf{THz$^2$} for dynamical
matrices and self-energies.  Force constants are mass-weighted
upstream by
\texttt{phonon\_inputs.fc3\_compression.build\_fc3\_target} (or by
\texttt{gamma\_project\_M\_blocks} for the $\Gamma$-only path) and the
dynamical-matrix entries are scaled by the
% REVIEW: math chars (m_u, ^2, (2pi)^2) cannot sit in \texttt{} text mode -- wrap the formula in math
\texttt{CONVERSION\_THZ2}~$= e/(m_u\,\text{\AA}^2)/(2\pi)^2\cdot 10^{-24}$
factor from \texttt{phonon\_inputs.constants}.  Self-energies always
carry a $q$-axis of shape
$(n_\mathrm{slabs}, n_\mathrm{kpts}, n_\mathrm{freq}, n_\mathrm{dof}, n_\mathrm{dof})$,
even when $n_\mathrm{kpts}=1$ (the $\Gamma$-only path).  This
unification means the SCBA fixed-point loop is a single function
shared by both drivers; the $\Gamma$-only path is recovered as
\texttt{transmission\_q(q\_mesh=(1,1,1))}, and the regression check
\texttt{compare\_q11\_to\_finite} asserts that the two paths produce
identical $T(\omega)$.

The frequency grid is the symmetric
$[-f_\mathrm{max}, \ldots, -\Delta\omega, 0, \Delta\omega, \ldots,
+f_\mathrm{max}]$ array built by
\texttt{grids.build\_frequency\_grid(freq\_range\_thz, eta\_w\_thz,
eta\_factor)}.  The $\omega = 0$ sample sits at index
$\mathrm{mid} = n_\mathrm{freq\_pos}$ and is included to make FFT
arithmetic clean; physical integrals exclude it via the
\texttt{pos\_mask} returned by the same routine, and the
bubble integrand zeroes the $\omega=0$ slice of the input Green's
functions before FFT to suppress the bosonic singularity.  The
% REVIEW: this appendix uses a LINEAR eta=eta_factor*dOmega convention; the theory chapter (eq. ~40) uses a QUADRATIC eta=eta_factor*dOmega^2 -- different convention noted explicitly
% REVIEW: default eta_factor=1.0 (so eta~dOmega); 0.05 is the LEGACY value, not the working default (consistent with the options table)
Lorentzian broadening parameter $\eta$ defaults to
\texttt{eta\_factor}$\cdot \Delta\omega$ (linear in the grid spacing,
with the default \texttt{eta\_factor}$=1.0$ giving $\eta\sim\Delta\omega$;
the legacy value was $\eta = 0.05\,\Delta\omega$), but can be overridden
directly via \texttt{eta\_w\_thz}.  Note that the theory chapter adopts a
different, \emph{quadratic} convention $\eta = \texttt{eta\_factor}\cdot
(\Delta\omega)^2$; the appendix retains the linear form used in the solver
code.

\section{The bubble kernel: \texttt{bubble.py}}
\label{sec:bubble_kernel}

The shared FFT bubble kernel
\texttt{bubble\_dense(phi\_left, phi\_right, G\_a, G\_b, n\_fft,
prefactor, out\_slice, zero\_freq\_idx, xp)} computes one $(I, J)$ block
of the 3-phonon self-energy
\begin{equation}\label{eq:bubble_block}
  \Sigma^x_{IJ}(\omega) = \frac{i\hbar}{2}\sum_{c,d,e,f}
  \phi^L_{Ice}\,
  G^x_{ed}(\omega) \star G^x_{cf}(\omega)\,
  \phi^R_{Jdf},
\end{equation}
where $x \in \{<, >\}$ selects the lesser or greater branch and
$\star$ denotes the frequency convolution.  The kernel evaluates the
convolution as a zero-padded FFT along the $\omega$ axis -- the
% REVIEW: FLAG-ONLY -- bubble_dense rounds n_fft to the next power of two, but se_finite/theory (eq:nfft_pad) state n_fft = 2 n_freq - 1; unreconciled FFT length
FFT length \texttt{n\_fft} is the next power of two above the linear
output length $n_\mathrm{freq}^\mathrm{out} + n_\mathrm{freq}^\mathrm{G\_a} - 1$ --
contracts indices with two \texttt{einsum} calls per frequency
sample (one for $\phi^L \cdot G^x$, one for the result times $\phi^R
\cdot G^x$), and slices the output to the caller's
\texttt{out\_slice}.  The vector backend is selected via \texttt{xp}
(default \texttt{numpy}; a \texttt{cupy} backend gives drop-in GPU
support when available).

The prefactor multiplying the FFT result is supplied by the caller;
the canonical value is $0.5\,i\,\hbar\,d\omega / (2\pi)$ in SI units.
The kernel is byte-for-byte shared with the production block-sparse
solver via an independent installed copy at
\texttt{quatrex.phonon.bubble}.  Cross-agreement between the two
implementations is asserted by
\texttt{tests.quatrex.phonon.test\_bubble\_kernel}.

\section{Frequency grid and Bose distribution: \texttt{grids.py}}
\label{sec:solver_grids}

In addition to \texttt{build\_frequency\_grid}, this module exposes:

\begin{description}
  \item[\texttt{bose\_full\_axis(freqs\_thz, temperature\_k)}.] Returns
        $n_B(\omega) = 1/(e^{\hbar\omega/k_BT}-1)$ on the full
        symmetric grid.  The $\omega = 0$ value is set to $0$ (rather
        than infinity) so callers can multiply with $G^<$ at that
        sample without producing NaN; the caller is responsible for
        zeroing the $\omega=0$ contribution before frequency
        integration.
  \item[\texttt{boson\_contact\_self\_energies\_from\_gamma}.]  Builds
        $(\Sigma^<_\mathrm{lead}, \Sigma^>_\mathrm{lead})$ from the
        contact $\Gamma_\mathrm{lead}(\omega)$ matrix via the
        equilibrium fluctuation--dissipation theorem,
        $\Sigma^<_\mathrm{lead}(\omega) = i\,n_B(\omega)\,\Gamma_\mathrm{lead}(\omega)$
        and
        $\Sigma^>_\mathrm{lead}(\omega) = i\,(n_B(\omega)+1)\,\Gamma_\mathrm{lead}(\omega)$,
        with the lead temperature passed in.
\end{description}

\section{Retarded self-energy reconstruction: \texttt{retarded.py}}
\label{sec:retarded}

The retarded scattering self-energy is reconstructed from the
$(\Sigma^<, \Sigma^>)$ pair via the Kramers--Kronig relation
\begin{equation}\label{eq:retarded_kk}
  \Sigma^R(\omega)
  = \tfrac{1}{2}\bigl[\Sigma^>(\omega) - \Sigma^<(\omega)\bigr]
  + \frac{i}{2\pi}\,\mathrm{PV}\!\!\int\frac{\Sigma^>(\omega')-\Sigma^<(\omega')}{\omega-\omega'}\,d\omega'.
\end{equation}
\texttt{retarded.build\_retarded(sigma\_lesser, sigma\_greater,
omega\_grid\_thz, method)} dispatches over three options:

\begin{description}
  \item[\texttt{"half"}.]  Drops the principal-value integral entirely
        and returns $\Sigma^R = \tfrac{1}{2}(\Sigma^> - \Sigma^<)$.
        This is the imaginary part of $\Sigma^R$ only -- i.e.\ the
        % REVIEW: broadening ~ (Sigma^> - Sigma^<), not the sum (sign fix)
        broadening, $-\mathrm{Im}\,\Sigma^R = \tfrac{1}{2}(\Sigma^> - \Sigma^<)$
        up to a sign and factor.  Used when only the broadening matters
        (e.g.\ the $\Gamma$-only regression check
        \texttt{compare\_q11\_to\_finite}, where $\mathrm{Re}\,\Sigma^R$
        cancels in the trace).  Both
        \texttt{transmission\_finite} and \texttt{transmission\_q}
        default to this choice so the regression assertion is exact.
  \item[\texttt{"pv"}.]  Singularity-subtracted principal-value
        integral.  Rewrites the integrand as
        $[\Delta(\omega') - \Delta(\omega)]/(\omega-\omega')$
        (regular at $\omega'=\omega$) plus the pole-free remainder
        $\Delta(\omega)\,\mathrm{PV}\!\int 1/(\omega-\omega')d\omega'$.
        Costs $O(n_\mathrm{freq}^2)$ per matrix element.  Captures
        both $\mathrm{Re}\,\Sigma^R$ and $\mathrm{Im}\,\Sigma^R$.
  \item[\texttt{"fft"}.]  FFT Hilbert transform via the
        sign-multiplier method, $O(n_\mathrm{freq}\log n_\mathrm{freq})$
        per matrix element.  Preferred for large frequency grids; the
        sign of $-1j \cdot \mathrm{sign}(\omega)$ multiplier is set so the
        result matches the principal-value convention to $\sim 10^{-10}$
        in floating-point.
\end{description}

\noindent
The factor multiplying the
$(n_B/(n_B+1))$ sum in the eigenmode expansion of $\Sigma^R$ is
$-2i$ (no extra $\pi$); the $1/\pi$ from
$\mathrm{Im}\,G^R = -\pi\sum_n\delta(\omega-\omega_n)$ is absorbed
into the Lorentzian
$L_\eta(\omega-\omega_n) = (\eta/\pi)/[(\omega-\omega_n)^2+\eta^2]$
that replaces the delta in the broadened expansion.  This convention
% REVIEW: FLAG-ONLY -- hard-coded equation numbers 873-874 and file name theory.tex are fragile; use \cref labels instead
is the one used throughout \texttt{theory.tex} (see equations 873--874
of the theory chapter).

\section{Open-boundary leads and device Dyson equation: \texttt{leads.py}}
\label{sec:solver_leads}

This module owns the lead/contact construction and the
device-side Dyson solve.  The five public entry points:

\begin{description}
  \item[\texttt{sancho\_rubio(z2, H\_00, H\_01)}.]  Surface Green's
        function for one lead at frequency $z = \omega^2$, computed
        by iterative decimation \cite{sancho_highly_1985}.  The fast iteration
        is the standard $\varepsilon, \varepsilon_s, \alpha, \beta$
        update; convergence is monitored by the relative residual
        $\|\alpha\| / \|\varepsilon\|$.  If the matrix iteration
        diverges (residual $> 10^{-4}$ after \texttt{max\_iter}
        iterations), the routine raises \texttt{RuntimeError} so the
        caller can fall back to the undressed contact and emit a
        warning.

  \item[\texttt{sancho\_rubio\_batch(z2\_array, H\_00, H\_01)}.]
        Frequency-batched version; vectorises the matrix operations
        across the input $\omega^2$ array using \texttt{numpy.linalg}
        broadcasting.

  \item[\texttt{build\_device\_hamiltonian(H\_00, H\_01, n\_slabs)}.]
        Assembles a block-tridiagonal device Hamiltonian of
        \texttt{n\_slabs} identical slabs with on-site $H_{00}$ and
        first-neighbour coupling $H_{01}$.  Used in the
        $\Gamma$-only path where the device extends the same slab
        many times along the transport axis.

  \item[\texttt{compute\_obc\_batch(omega\_grid\_thz, H\_00, H\_01,
        temperature\_k, ...)}.]  On a frequency batch, computes the
        lead retarded self-energies $\Sigma^R_L$, $\Sigma^R_R$ via
        Sancho--Rubio plus contact coupling, then the lead
        broadenings $\Gamma_{L,R} = i(\Sigma^R_{L,R} - \Sigma^A_{L,R})$,
        and finally the lead lesser/greater pair via the
        fluctuation--dissipation relation from
        \texttt{boson\_contact\_self\_energies\_from\_gamma}.  An
        optional \texttt{lead\_sigma\_r\_L} /
        \texttt{lead\_sigma\_r\_R} input dresses the surface block
        with a scattering retarded self-energy from the SCBA loop; if
        Sancho--Rubio fails for the dressed lead the routine falls
        back to the undressed contact and emits a warning.

  \item[\texttt{solve\_green\_functions} /
        \texttt{solve\_green\_batch}.]  Given the total retarded
        self-energy
        $\Sigma^R_\mathrm{tot} = \Sigma^R_L + \Sigma^R_R + \Sigma^R_\mathrm{ph\text{-}ph}$
        and the total lesser/greater self-energies, solves the device
        Dyson equation
        $G^R(\omega) = [(\omega+i\eta)^2 - H_\mathrm{dev} - \Sigma^R_\mathrm{tot}]^{-1}$
        and the Keldysh equation
        $G^{<,>}(\omega) = G^R \Sigma^{<,>}_\mathrm{tot} G^A$.

  \item[\texttt{ballistic\_transmission\_z2(z2\_array, ...)}.]  Caroli
        formula transmission $T(\omega^2) = \mathrm{Tr}(\Gamma_L G^R \Gamma_R G^A)$
        on the input $\omega^2$ grid; serves as the ballistic
        baseline and the regression target for the SCBA loop.
\end{description}

% REVIEW: wrap heading math in \texorpdfstring for PDF bookmarks/intcalc
\section[Gamma-only self-energy driver: \texttt{se\_finite.py}]{Γ-only self-energy driver: \texttt{se\_finite.py}}
\label{sec:se_finite}

\texttt{compute\_phph\_self\_energy\_finite(G\_lesser, G\_greater, Phi,
omega\_grid\_thz, dw\_thz)} is a thin wrapper around
\texttt{bubble\_dense} encoding the $\Gamma$-only convention:

\begin{itemize}
  \item Symmetric $\omega$ grid of length $n_\mathrm{freq}$,
        constructed by \texttt{grids.build\_frequency\_grid}.
  \item FFT length $n_\mathrm{fft} = 2\,n_\mathrm{freq} - 1$ for the
        full causal convolution.
  \item $\omega = 0$ sample zeroed before FFT (singular Bose at
        $\omega=0$) via \texttt{zero\_freq\_idx} passed into the
        kernel.
  \item Output slice $[\mathrm{mid}\,:\,\mathrm{mid}+n_\mathrm{freq}]$
        with $\mathrm{mid} = n_\mathrm{freq}//2$ to recover the symmetric
        output grid.
\end{itemize}

Returns $(\Sigma^<, \Sigma^>)$ only -- the retarded reconstruction is
done downstream by \texttt{retarded.build\_retarded} so that callers
can pick the Kramers--Kronig precision they need.

% REVIEW: wrap heading math in \texorpdfstring for PDF bookmarks/intcalc
\section[q-resolved self-energy driver: \texttt{se\_q.py}]{$q$-resolved self-energy driver: \texttt{se\_q.py}}
\label{sec:se_q}

The $q$-resolved path generalises \texttt{se\_finite.py} by lifting the
real-space FC3 to a Fourier-space vertex
$\Phi(\vc{q}, \vc{q}', \vc{q}'')$ on a $\Gamma$-centred
$q$-mesh of size $n_\mathrm{kpts}$.  Conservation of momentum imposes
$\vc{q}'' = -\vc{q} - \vc{q}' \pmod{\vc{G}}$ so only two of the three
indices are free; the worker enumerates the free pair and computes
the bubble contribution.

\texttt{compute\_phph\_self\_energy\_q\_dense(...)} dispatches the
computation as follows.  For each external $q$-point (an outer index
parallelised by the worker pool), the worker iterates over all
internal $q'$-points compatible with momentum conservation, picks the
corresponding $G^{<,>}(q', \omega)$ slices, contracts with the
pre-computed FFT-domain $\Phi$ tensor along the frequency axis, and
% REVIEW: FLAG-ONLY -- literal Unicode 'Sigma' glyph mixed with math $^{<,>}$ relies on unicode-math; prefer $\Sigma^{<,>}$
accumulates into the output Σ$^{<,>}(q_\mathrm{ext}, \omega)$.  The
$q$-difference index map \texttt{q\_diff\_map} is pre-built so the
inner loop is a flat array access rather than a modular-arithmetic
lookup.

The dispatcher picks between a \texttt{ThreadPoolExecutor} (preferred
for memory-bound workloads when the FFT batch is large) and a
\texttt{multiprocessing.get\_context('forkserver')} pool (preferred
for CPU-bound workloads when each worker can amortise the
NumPy thread-pool startup over many external $q$-points).  The choice
is driven by the environment variable \texttt{PHONON\_SOLVER\_POOL}
(\texttt{thread}/\texttt{process}, default \texttt{thread}).

\section{Top-level drivers: \texttt{dense.py}}
\label{sec:dense_driver}

\texttt{dense.py} is the user-facing module that exposes the
$\Gamma$-only and $q$-resolved drivers, the shared SCBA loop, and the
helpers that load and project the FC3 vertex.

\paragraph{\texttt{transmission\_finite(fc3\_source, fc2, masses,
freq\_range\_thz, temperature\_k, ...)}.}  The $\Gamma$-only driver.
Loads FC3 (via \texttt{load\_fc3\_raw}, which accepts an HDF5 path, a
numpy array, or a phono3py-style dict), $\Gamma$-projects it to the
primitive-cell DOFs via \texttt{gamma\_project\_M\_blocks}, builds the
device-side Hamiltonian and lead self-energies, and runs either one
shot of the bubble (\texttt{n\_iter=0}) or the SCBA loop
(\texttt{n\_iter $>$ 0}).  Returns
$(\omega, T(\omega), \Sigma^<, \Sigma^>, G^<, G^>, \Sigma^R_\mathrm{lead})$
plus a diagnostic dict.

\paragraph{\texttt{transmission\_q(..., q\_mesh, n\_qpoints, ...)}.}
The $q$-resolved driver.  Takes a $q$-mesh tuple (e.g.\
\texttt{(8, 8, 1)} for an 8$\times$8 transverse mesh with a single
sample along the transport axis) and runs the bubble on the
$q$-dependent Green's functions.  The lead self-energies are built
per $q$-point.  Returns the same output dictionary as
\texttt{transmission\_finite} with the $q$-axis added.

\paragraph{\texttt{scba\_loop(...)}.}  The SCBA fixed-point iteration
shared by both drivers.  At each step:
\begin{enumerate}
  \item Build $\Sigma^R$ from the current $(\Sigma^<, \Sigma^>)$ via
        the user-selected reconstruction method.
  \item Add the lead retarded self-energy and solve the device Dyson
        equation for $G^R$.
  \item Solve the Keldysh equation for $(G^<, G^>)$.
  \item Symmetrise $(G^<, G^>)$ via
        \texttt{diagnostics.symmetrize\_lesser\_greater} to project
        out the components that break the bosonic Keldysh symmetry.
  \item (Optional, \texttt{zero\_mode\_projection=True}) project
        $\Sigma^{<,>}$ onto the translation-free subspace.
  \item Recompute $(\Sigma^<, \Sigma^>)$ from the bubble on the new
        $(G^<, G^>)$.
  \item Apply the user-selected accelerator
        (\texttt{solver}: \texttt{"linear"}, \texttt{"anderson"},
        \texttt{"jfnk"}, or \texttt{"anderson+jfnk"}) and check
        convergence on the SCF residual
        $\|F(\Sigma)\|/\|\Sigma\|$.
\end{enumerate}
The convergence/stop test, the divergence guard, the Anderson
safeguards, and the JFNK fallback are described in
\Cref{ssec:solver_safeguards}. The drivers also auto-extend
$f_{\max}$ before sizing the frequency grid to ensure
$f_{\max}\ge 2\omega_{\max}/2\pi$ holds for the device dynamical
matrix (\Cref{ssec:solver_convergence}). The loop emits a warning
if it exceeds \texttt{max\_iter} without converging and returns the
best iterate (lowest residual seen), not the last.

\paragraph{\texttt{gamma\_project\_M\_blocks(fc3\_raw, masses,
nat\_prim)}.}  Direct $\Gamma$-only supercell$\to$primitive
projection of the FC3 vertex.  Sums the supercell FC3 over the
transverse lattice index of the second and third legs, then
mass-weights to produce the rank-3 lifted target of shape
$(n_\mathrm{dof}^\mathrm{prim}, \dim_{sc}, \dim_{sc})$ in THz$^{5/2}$.
This is the bubble's $\phi^L = \phi^R$ vertex for the $\Gamma$-only
path.

\paragraph{\texttt{compare\_q11\_to\_finite(...)}.}  Regression check
asserting \texttt{transmission\_q(q\_mesh=(1,1,1))} reproduces
\texttt{transmission\_finite(...)} to within numerical roundoff (the
working tolerance is $10^{-8}$ on $T(\omega)$).  Run by the test
suite and by the user when they touch any of the shared code.

\section{Diagnostics and physical-invariant checks: \texttt{diagnostics.py}}
\label{sec:solver_diagnostics}

\texttt{diagnostics.py} hosts the runtime checks the SCBA loop relies
on but does not enforce:

\begin{description}
  \item[\texttt{check\_broadening\_sign(Sigma\_R, freqs\_thz, name)}.]
        For $\omega > 0$ the broadening matrix
        $\Gamma = i(\Sigma^R - \Sigma^A)$ must be positive semi-definite
        (PSD), and for $\omega < 0$ negative semi-definite (NSD), as
        required by causality of the retarded self-energy.  The
        routine diagonalises $\Gamma$ at each frequency and reports
        the number of sign violations and the worst-case eigenvalue.
  \item[\texttt{check\_full\_axis\_symmetry(G\_lesser, G\_greater)}.]
        Asserts the bosonic Keldysh symmetry
        $G^<(\omega) = [G^>(-\omega)]^T$ and
        $G^R(\omega) = [G^R(-\omega)]^*$ to within tolerance.  Used
        as a post-SCBA sanity check.
  \item[\texttt{symmetrize\_lesser\_greater(G\_lesser, G\_greater)}.]
        Projection onto the full-axis symmetry manifold.  Symmetrises
        the lesser/greater pair by averaging $G^<$ and $[G^>(-\omega)]^T$
        and similarly for $G^>$.  Called after each SCBA iteration to
        prevent symmetry drift accumulating with iteration count.
\end{description}

\section{Cutoff policies: \texttt{cutoffs.py}}
\label{sec:solver_cutoffs}

\texttt{cutoffs.py} defines the cutoff policy that the cutoff-sweep
harness in \texttt{finite\_analysis.sse\_cutoffs} feeds back into the
bubble kernel.  The unified solver defaults to \emph{no} approximation:
full off-diagonal $G$ blocks inside the bubble, full FC3 vertex, no
magnitude or distance threshold.  Four knobs are exposed:

\begin{description}
  \item[\texttt{diag\_G\_in\_se}.]  Restrict the inner $G$ in the bubble
        integrand to the diagonal $(K_1=K_1', K_2=K_2')$ blocks.
        Recovers the original block-tridiagonal NEGF approximation
        when \texttt{True}.
  \item[\texttt{fc3\_nn\_only}.]  Drop FC3 block triplets with
        $|I-J| > 1$, $|I-K| > 1$, or $|J-K| > 1$.  Already the default
        of the production \texttt{quatrex.phonon.fc3\_loader.fc3\_to\_phi\_blocks}.
  \item[\texttt{fc3\_distance\_cutoff}.]  Drop FC3 entries whose
        triplet diameter (the max of the three pairwise atomic
        distances) exceeds the cutoff in \AA{}.
  \item[\texttt{fc3\_magnitude\_threshold}.]  Drop FC3 entries whose
        magnitude is below
        \texttt{threshold}~$\times \max|\Phi^3|$.
\end{description}

Each knob is independent; the harness in
\texttt{finite\_analysis.sse\_sparsity\_driver.run\_cutoffs} sweeps
them and tabulates the effect on $T(\omega)$ and the Landauer heat
current $Q$, answering the question of which approximations preserve
the transport observables.

\section{Public API surface}
\label{sec:solver_api}

The exports in \texttt{phonon/solver/\_\_init\_\_.py} bundle the
above:

\begin{verbatim}
from phonon.solver import (
    transmission_finite,    # Γ-only SCBA driver
    transmission_q,         # q-resolved SCBA driver
    scba_loop,              # shared fixed-point loop
    gamma_project_M_blocks, # Γ-only supercell -> primitive
    compare_q11_to_finite,  # regression check
    bubble_dense,           # FFT 3-phonon bubble kernel
    build_retarded,         # Σ^R reconstruction (half|pv|fft)
    sancho_rubio, sancho_rubio_batch, compute_obc_batch,
    build_frequency_grid, bose_full_axis,
)
\end{verbatim}

The companion production solver
\texttt{quatrex.phonon.sse\_phonon\_phonon} is the MPI-distributed
block-sparse variant and consumes the same bubble kernel through its
own installed copy at \texttt{quatrex.phonon.bubble}.  It realises the
full off-diagonal ring through the FFT-first energy pipeline and the
1-D spatial halo whose cost model is derived in
\cref{ssec:phph_dist_impl}, and is pinned \emph{bit-identical} to this
dense reference both serially and across \texttt{comm.stack} /
\texttt{comm.block} factorisations (\cref{tab:phph_dist}).  The dense
solver of this appendix remains the arbiter of correctness for the
production stack.

% ======================================================================
\section{Physics verification of the dense solver}
\label{sec:phonon_solver_verification}
% ======================================================================

Because the dense solver is the arbiter of correctness, its
3-phonon physics, its controllable approximations, its convergence
behaviour and its zero-mode handling were audited end to end. The
audit is reproducible: it is implemented as six standalone scripts
under \texttt{phonon/scripts/} (\texttt{verify\_bubble.py},
\texttt{verify\_cutoffs.py}, \texttt{verify\_scba\_convergence.py},
\texttt{verify\_discretization.py}, \texttt{verify\_zero\_modes.py},
\texttt{verify\_literature.py}), each of which prints a pass/fail
report and emits diagnostic plots to \texttt{phonon/scripts/out/verify/},
and as a permanent regression suite in
\texttt{tests/phonon\_inputs/test\_solver\_physics.py}. The analytic
toy systems that serve as ground truth (single anharmonic oscillator,
monatomic and diatomic chains) live in
\texttt{phonon/solver/toy\_models.py}.

\subsection{Self-energy correctness and absence of double counting}

The 3-phonon bubble was checked against an independent brute-force
frequency convolution: the zero-padded FFT kernel reproduces the
explicit sum to floating-point precision, and a delta-function probe
confirms exact frequency conservation
($\omega_\mathrm{out}=\omega_a+\omega_b$). The prefactor
$\tfrac{i\hbar}{2}\cdot\tfrac{d\omega}{2\pi}$ is the Migdal-bubble
form $\Sigma^{\lessgtr}_{ab}=\tfrac{i\hbar}{2}\sum_{c,d,e,f}
\Phi_{acd}[G^{\lessgtr}\!*\!G^{\lessgtr}]_{cf,de}\Phi_{bef}$:
the factor $\tfrac12$ is the loop
symmetry factor of the unique two-vertex diagram (a Wick count gives
$(1/3!)^2\cdot 3\cdot 3\cdot 2 = 1/2$), and
\texttt{build\_realspace\_fc3\_matrices} stores $\Phi$ as the bare
mass-weighted third derivative with no combinatorial factor baked in,
so there is no double counting.

The strongest no-double-counting check is detailed balance: with
Green's functions in thermal equilibrium the bubble must return
$\Sigma^{>}(\omega)=e^{\hbar\omega/k_BT}\,\Sigma^{<}(\omega)$. This
holds to $\sim10^{-9}$ relative (FFT round-off on sharply peaked
input), which certifies that the Bose occupation enters exactly once
and that the $\Sigma^{<}/\Sigma^{>}$ signs are correct. The raw
bubble output is already Keldysh-symmetric
($\Sigma^{<}(\omega)=[\Sigma^{>}(-\omega)]^{T}$) and anti-Hermitian to
machine precision, so the per-iteration
\texttt{symmetrize\_lesser\_greater} projection is a genuine no-op on
a consistent input rather than a band-aid hiding an error.

\subsection{The FFT Hilbert transform: periodic-wrap correction}

The retarded self-energy is reconstructed from $\Sigma^{<,>}$ by the
Kramers--Kronig relation
$\Sigma^{R}=\tfrac12(\Sigma^{>}-\Sigma^{<})
+\tfrac{i}{2}\mathcal{H}[\Sigma^{>}-\Sigma^{<}]$. The
default \texttt{retarded="fft"} path used the bare sign-multiplier
FFT, which computes the \emph{periodic} Hilbert transform. For the
finite-support self-energies fed here the periodic wrap injects a
spurious discontinuity and a resolution-independent ${\sim}1\%$ error
in $\mathrm{Re}\,\Sigma^{R}$. The fix
(\texttt{hilbert\_transform\_axis} in \texttt{solver/retarded.py})
zero-pads the signal before the FFT, turning the circular convolution
into a truncated linear one; with the default eightfold padding the
interior error drops to ${\sim}10^{-5}$, matching the
$\mathcal{O}(n_\omega^2)$ principal-value reference. The two
reconstruction methods then agree to $\sim0.4\%$ on a resolved grid.

\subsection{Controllable approximations}

Every approximation against the strict SCBA expression has an explicit
control, and each control's ``off'' setting was shown to recover the
unapproximated result. \texttt{sigma\_cutoff}, \texttt{vertex\_cutoff}
and \texttt{g\_cutoff} all default to \texttt{None}; with all cutoffs
disabled the multi-slab kernel reproduces an independent brute-force
sum over every block quadruple to floating-point precision.
\texttt{sigma\_cutoff} only filters the output $(I,J)$ set --- retained
blocks are bit-identical to the untruncated computation;
\texttt{g\_cutoff} saturates to the full result once it reaches the
device size; \texttt{vertex\_cutoff} only drops FC3 triplets.
\texttt{dc\_handling} and \texttt{retarded} select distinct, bounded
results. The $\Gamma$-only treatment of transverse momentum is not an
approximation for a nanowire in vacuum (there is no transverse
periodicity); the non-self-consistent leads remain the one standing
approximation, controlled by \texttt{scattering\_contacts}.

\subsection{Convergence}
\label{ssec:solver_convergence}

On the analytic toy systems the SCBA fixed-point loop behaves as
expected: it converges for weak-to-moderate anharmonicity, $\Sigma^{R}$
at loop exit is consistent with the stored $\Sigma^{<,>}$ (no lag), and
the per-iteration symmetrization is a genuine no-op. The real d5a SiNW
is a much harder case; running it --- rather than trusting the toys
--- exposed two distinct problems whose joint effect was a divergent
multi-slab SCBA.

\paragraph{Single slab: the mixing scheme.}
For the single-cell d5a device with a properly-sized grid (see below),
the historical default linear mixing $\alpha=0.5$ \emph{does not}
converge: $d\Sigma/\Sigma$ oscillates between $\sim15\%$ and $\sim67\%$
from one iteration to the next. The bare (unregularised) Anderson
scheme is worse --- it diverges, with the heat-flow conservation
error blowing up to $\sim100\%$, because its unregularised least-squares
extrapolates an overshooting step from a residual of large dynamic
range. Damping the linear mix to $\alpha=0.3$ (the new default)
suppresses the oscillation, converging the single-slab SCBA in
${\sim}65$ iterations; a safeguarded Anderson scheme (described in
\Cref{ssec:solver_safeguards}) converges it in ${\sim}47$.

\paragraph{Multi-slab: the $f_{\max}$ aliasing.}
For $n_\mathrm{slabs}\ge 2$ even the corrected mixing did not converge
on d5a --- the residual reached $\sim10^{-4}$, then \emph{grew}, with
$|\Sigma^{R}|$ climbing from a few hundred to many thousands of
$\mathrm{THz}^{2}$ over $\sim 20$ iterations. The same symptoms appear
with safeguarded Anderson, with Jacobian-free
Newton--Krylov, with every cutoff combination, and at every mixing
factor --- a hallmark of a \emph{linearly unstable} fixed point
(spectral radius of the SCBA Jacobian $>1$), not of a tuning problem.
The diagnosis tracks to the discretisation requirement
$f_{\max}\ge 2\omega_{\max}/2\pi$ derived in
\cref{ssec:phph_freq}: the historical d5a configurations used
$f_{\max}\in\{16,18,24\}$\,THz, while
$2\omega_{\max}/2\pi\approx 129$\,THz for the Si--H stretch. The
bubble's convolution support \cref{eq:conv_support} sticks out of the
grid by a factor of $\sim 7$; the truncated and aliased part injects a
fake strong-coupling $\Sigma$ that breaks the SCBA contraction.

The bulk-Si control closes the diagnosis. Repeating the
one-shot (lowest-order Born) bubble through the identical pipeline ---
\texttt{transmission\_finite(max\_scba\_iter=1)} --- for bulk Si and
for d5a, with and without the necessary grid extension, gives the
absolute magnitudes summarised in
\cref{tab:fmax_aliasing_one_shot}. With the under-ranged grid, even
bulk Si --- a textbook weakly anharmonic crystal --- gives
$|\Sigma^{R}|/\omega_{\max}^{2}\approx 33\%$ and a median linewidth
$\sim 0.67$\,THz, an order of magnitude above the experimental
$\sim 0.05$--$0.1$\,THz for the Si optical mode at $300$\,K. Extending
% REVIEW: 0.024 THz is BELOW the 0.05-0.1 THz experimental range, not "within" it; one-shot Born underestimates, ~factor of two (consistent with caption)
$f_{\max}$ to cover the convolution support recovers
$\sim 0.024$\,THz, below the experimental $\sim 0.05$--$0.1$\,THz range
(the one-shot Born underestimates), within about a factor of two. For d5a the
correction is even larger, $\sim 25\times$ in $|\Sigma^{R}|$, and the
relative magnitude drops to $|\Sigma^{R}|/\omega_{\max}^{2}\approx
0.2\%$ --- squarely in the perturbative regime where SCBA is well
defined.

\begin{table}[ht]
  \centering
  \begin{tabular}{lrrrr}
    \hline
    system & $f_{\max}$ (THz) & $2\omega_{\max}/2\pi$ (THz)
      & $|\Sigma^{R}|_{\max}$ (THz$^{2}$)
      % REVIEW: FLAG-ONLY -- |Im Sigma^R|/omega is not the linewidth; HWHM = |Im Sigma^R|/(2 omega) (factor-2 ambiguity)
      & median $|\mathrm{Im}\,\Sigma^{R}|/\omega$ (THz) \\
    \hline
    Bulk Si & $18$ (under-ranged) & $30$  & $72$  & $0.67$ \\
    Bulk Si & $36$ (proper)       & $30$  & $22$  & $0.024$ \\
    d5a SiNW & $18$ (under-ranged) & $129$ & $205$ & $1.14$ \\
    d5a SiNW & $140$ (proper)      & $129$ & $8.7$ & $1.1\times10^{-5}$ \\
    \hline
  \end{tabular}
  % REVIEW: FLAG-ONLY -- "~25-30x larger Sigma" over-generalises: table shows bulk Si only ~3.3x (72->22); only d5a is ~24x
  \caption{One-shot (lowest-order Born) self-energy magnitude through
  the identical \texttt{transmission\_finite} pipeline, at
  $n_\mathrm{slabs}=1$ and $\eta/\Delta\omega=1$. Under-ranged grids
  fake a $\sim25$--$30\times$ larger $\Sigma$ that destabilises SCBA;
  % REVIEW: Debernardi 1995 is a first-principles DFPT lifetimes paper, not experimental
  the bulk-Si linewidth under the proper grid matches the first-principles (DFPT)
  $\sim 0.05$\,THz \cite{debernardi_PRB_1995} within a factor of two
  (the one-shot Born underestimates the converged value).}
  \label{tab:fmax_aliasing_one_shot}
\end{table}

\paragraph{Safety net: auto-extending $f_{\max}$.}
With the diagnosis confirmed --- the FC3 generation, the
\texttt{eV/(\AA{}\textsuperscript{3}\,amu\textsuperscript{3/2})}-to-$\mathrm{THz}^{2}$
conversion in \texttt{build\_realspace\_fc3\_matrices}, and the bubble
prefactor are all unit-correct, only the grid was wrong --- the
solver entry points \texttt{transmission\_finite} and
\texttt{transmission\_q} were made safe by default. They compute
$\omega_{\max}$ from $H_{00}+H_{01}+H_{01}^{\dagger}$ before sizing the
grid and, if the caller's $f_{\max}$ is below $1.05\times
2\omega_{\max}/2\pi$, transparently extend the range and rescale
$N_\omega$ to preserve $\Delta\omega$. The behaviour is gated by
\texttt{auto\_extend\_fmax=True} (default) and a \texttt{fmax\_margin}
parameter; passing \texttt{auto\_extend\_fmax=False} reverts to a loud
\texttt{RuntimeWarning} that names the required minimum. The helper
\texttt{\_ensure\_fmax} centralises the logic in one place.

\subsection{Solver safeguards for strongly anharmonic systems}
\label{ssec:solver_safeguards}

Even with the grid sized correctly, the d5a multi-slab problem is
strongly anharmonic ($|\Sigma|/|H|\sim 0.05$--$0.1$, not the few
per~cent of a bulk crystal), and a plain Picard / Anderson scheme can
still struggle. \texttt{dense.py} adds five orthogonal toggles
covering the regime; setting them all to their legacy values
reproduces the earlier solver bit-for-bit.

\paragraph{Safeguarded Anderson acceleration.}
The \texttt{\_AndersonAccelerator} class replaces the earlier
hard-restart \texttt{\_anderson\_mix}, which cleared the entire history
on \emph{any} residual uptick. On the noisy residual near the
fixed point the restart fired every iteration, collapsing Anderson to
a bare linear step that --- on an unstable fixed point --- diverges.
The new scheme follows the standard
recipe~\cite{walker_anderson_2011,toth_convergence_2015}: history is
\emph{kept} across small upticks; the least-squares is solved by a
truncated-SVD \texttt{lstsq} with cutoff $\mathrm{rcond}=10^{-8}$,
filtering ill-conditioned secant directions instead of discarding all
history; an over-long extrapolation step (more than $2\times$ the
damped linear step) is replaced by the linear step for that iteration
only; and the history is reset only on \emph{genuine stagnation},
defined as five consecutive iterations without a new best residual.
The residual is measured in a block-weighted norm so that small
far-off-diagonal $\Sigma_{IJ}$ blocks --- invisible to a global 2-norm
--- still contribute to both the least-squares and the restart test.

\paragraph{Zero-mode projection of the self-energy.}
The 3-phonon bubble is not guaranteed to satisfy a discrete acoustic
sum rule on $\Sigma$: $\mathrm{Re}\,\Sigma^{R}$ generally has a
nonzero projection onto every near-zero mode of the cell dynamical
matrix at $\Gamma$. Added to $H_{D}$ in the Dyson denominator this
projection shifts the affected $\omega^{2}$ --- possibly negative,
which puts a $G^{R}$ pole on the real axis. A more insidious failure
mode for an H-terminated wire is the IR side: at $T=300$\,K the Bose
factor $n_{B}(\omega)\sim k_{B}T/(\hbar\omega)$ on a near-zero mode is
enormous, so the Bose-enhanced $G^{<}\propto n_{B}A$ carries an IR
singularity that the bubble integrates and dumps back into $\Sigma$,
making the SCBA map non-contractive even with the grid sized
correctly.

For a 1\,D wire the periodic cell dynamical matrix
$H_{00}+H_{01}+H_{01}^{\dagger}$ has \emph{four} near-zero eigenmodes
at $\Gamma$, not three: the three rigid translations (machine-zero
$\omega^{2}\sim 10^{-13}$ for an ASR-clean FC2) plus the rigid
rotation about the wire axis --- the ``twist'' mode --- at the
finite-but-tiny value set by the wire's torsional restoring force.
For the d5a SiNW the twist sits at
$\omega^{2}\approx 7.7\times10^{-4}\,\mathrm{THz}^{2}$, i.e.\
% REVIEW: n_B = 1/(exp(hf/kT)-1) ~= 223 for f=0.028 THz at 300 K (was ~700)
$\omega\approx 0.028$\,THz, which gives $n_{B}\approx 223$ at room
temperature --- the dominant IR contribution to the bubble.

The module \texttt{phonon/solver/zero\_modes.py} therefore builds the
projector from the cell dynamical matrix itself, picking every
eigenvector below a relative-cutoff threshold:
\begin{equation}
  Q \;=\; \mathbf{1} \;-\; \sum_{k:\,\omega_{k}^{2}<\epsilon\,\omega_{\max}^{2}}
  \hat{v}_{k}\,\hat{v}_{k}^{\top},
\end{equation}
with $\hat{v}_{k}$ the cell eigenvector with $\omega_{k}^{2}<\epsilon
\omega_{\max}^{2}$ replicated uniformly across the $n_\mathrm{slabs}$
device slabs (the $q=0$ phase) and then orthonormalised. The default
relative threshold is $\epsilon=10^{-4}$, matching
\texttt{verify\_zero\_modes.check\_acoustic\_modes}: for d5a this
captures the three translations and the twist, four modes in total.
The mass-weighted construction $t_{\alpha}[\,l,i,\beta\,]=\sqrt{m_{i}}\,
\delta_{\alpha\beta}$ used by \texttt{build\_translation\_projector}
is preserved as a fallback for harnesses that build $H_{D}$ directly
without a periodic cell.

The SCBA loop applies $\Sigma^{\gtrless}\leftarrow Q\Sigma^{\gtrless}Q$
after symmetrisation. The projection is linear, so $\Sigma^{R}$ rebuilt
by \texttt{build\_retarded} is automatically zero-mode-clean and
satisfies $\hat{v}_{k}^{\top}\Sigma^{R} \hat{v}_{k}=0$ identically.
Gated by \texttt{zero\_mode\_projection=True}.

\paragraph{Jacobian-free Newton--Krylov fallback.}
For genuinely linearly-unstable fixed points the Picard map cannot be
made contractive by any mixing factor, so a Newton-type step is
needed. The per-iteration body of \texttt{scba\_loop\_dev} is
refactored into a pure map
$F(\mathbf{x})=\mathtt{step}(\mathbf{x})-\mathbf{x}$ on the flattened
$[\Sigma^{<},\Sigma^{>}]$ vector. The \texttt{solver} parameter
selects between
\texttt{"linear"} / \texttt{"anderson"} / \texttt{"jfnk"} /
\texttt{"anderson+jfnk"}; the JFNK path calls
\texttt{scipy.optimize.newton\_krylov} on $F$, packing complex
self-energies as a real vector for the inner LGMRES solve.
\texttt{"anderson+jfnk"} runs safeguarded Anderson first and, if it
hits the divergence guard or exhausts \texttt{max\_scba\_iter},
restarts JFNK warm-started from the best Anderson iterate.

\paragraph{Convergence test decoupled from conservation.}
The historical loop only \texttt{break}-ed when the SCF residual,
$dJ/J$ \emph{and} the heat-flow conservation error were all below
their respective tolerances. Conservation error, however, is set
predominantly by the discretisation --- grid spacing, $\eta/\Delta\omega$,
$f_{\max}$ --- and cannot be reduced by further SCBA iterations.
Coupling it to the stop test caused the loop to iterate a
numerically-converged self-energy past its useful life and, on a
near-unstable fixed point, into divergence. The
\texttt{gate\_on\_conservation} flag preserves the legacy behaviour;
the new default stops on the SCF residual alone and reports the
conservation error as a diagnostic.

\paragraph{Causality enforcement of $\Sigma^{R}$.}
Even with the grid sized for the convolution support and the cell
zero modes projected out, the discretised bubble plus the
per-iteration symmetrisation and zero-mode projection can leak
negative eigenvalues into $\Gamma_{\Sigma}=i(\Sigma^{R}-\Sigma^{A})$
for $\omega>0$, violating retarded causality. Once $\Sigma^{R}$ is
acausal the Dyson solve produces a $G^{R}$ with poles in the upper
half-plane; the next bubble amplifies the violation, the
\emph{symptom} on d5a being a catastrophic
$\sim\!1000\times$ jump in the SCF residual after several iterations
of apparent improvement. A standard remedy
% REVIEW: original key pourfath_causal_2010 was unresolvable (no such 2010 paper);
% re-sourced to Pourfath's 2014 NEGF monograph, which treats causal/retarded
% self-energy enforcement. The PSD-eigenvalue clip + KK rebuild below is our
% specific implementation.
\cite{pourfath_negf_2014} is to restore causality of $\Sigma^{R}$ by projecting it onto the
causal manifold each iteration: write
$\Sigma^{R}=X+iY$ with $X,Y$ matrix-Hermitian, take the eigen-
decomposition of $Y(\omega)$, clip its positive eigenvalues
(\emph{i.e.}\ the negative eigenvalues of $\Gamma_{\Sigma}=-2Y$) to
zero for $\omega>0$ (and symmetrically for $\omega<0$), and rebuild
$X$ from the cleaned $Y$ via the Kramers--Kronig relation
$X=-(1/\pi)\mathcal{H}[Y]$ (the same zero-padded FFT kernel used by
\texttt{build\_retarded}). \texttt{phonon/solver/causality.py}
implements this as \texttt{enforce\_causality\_psd}; the option
\texttt{causality\_projection=True} applies it after every
\texttt{build\_retarded} call inside the SCBA loop. A free-standing
\texttt{causality\_diagnostic} reports the worst PSD violation and
the frequency where it lives, and
\texttt{dynamical\_stability\_diagnostic} reports the smallest
eigenvalue of $H_{D}+\mathrm{Re}\,\Sigma^{R}(\omega)$ across the grid
-- a value $<0$ signals that the self-consistent self-energy has
driven a mode's $\omega^{2}$ past zero and the next Dyson denominator
is singular.

\paragraph{Vertex-rescaling Pad\'e (LOA) baseline.}
For systems where the SCBA fixed point is genuinely linearly
unstable (the residual stays bounded away from zero under every
linear and Newton step), the principled alternative to iterating an
unstable map is the rescaling technique of
\cite{mera_inelastic_2012}. The bubble self-energy scales as
$\Sigma\propto\lambda^{2}$ in the cubic coupling, so the SCBA
fixed-point observable $\mathcal{O}(\lambda)$ can be expanded as a
series in $\lambda^{2}$. Running the SCBA at a few rescaled vertices
$\lambda\in(0,1]$ -- where small $\lambda$ converges easily because
$|\Sigma|/|H|\propto\lambda^{2}$ is small -- and Pad\'e-extrapolating
the result to $\lambda=1$ produces an estimate of the SCBA answer
without iterating an unstable fixed point. The option
\texttt{vertex\_scale} on \texttt{transmission\_finite} multiplies
the FC3 vertex by the requested $\lambda$; the diagnostic driver
\texttt{phonon/scripts/scba\_diagnostic.py} runs the sweep
$\lambda\in\{0.2,0.4,0.6,0.8\}$ and fits a quadratic-in-$\lambda^{2}$
Pad\'e to the anharmonic conductance, plotted under
\texttt{vertex\_pade.pdf}.

\paragraph{Cluster diagnostic driver.}
\texttt{phonon/scripts/scba\_diagnostic.py} runs the full
configuration matrix --- baseline, causality-enforced, linear at
$\alpha\in\{0.10,0.05,0.02\}$, pure JFNK, \texttt{anderson+jfnk},
lowest-order Born, and the Pad\'e sweep --- on the same system in one
invocation. Each run writes a per-iteration trajectory CSV
(residual, $|\Sigma^{R}|_{\max}$, conservation, causality violation,
$\min\lambda(H_{D}+\mathrm{Re}\,\Sigma^{R})$), a \texttt{result.npz}
with the observable arrays, and a summary JSON; the driver then
produces a four-panel overlay (\texttt{trajectory\_overlay.pdf})
showing how each method behaves across the same iterations. This is
the recommended workflow for any new hard system: run the driver,
inspect the overlay, and pick the configuration whose causality and
stability curves stay closest to zero.

\paragraph{Divergence guard and best-iterate return.}
The loop tracks the running minimum $\rho_{\min}$ of the residual and
aborts if $\rho>10\rho_{\min}$ for four consecutive iterations; the
returned $(\Sigma^{<},\Sigma^{>})$ is then the best iterate seen, not
the last (possibly NaN-poisoned) one. Together with the conservation
decoupling this means a bad run fails fast and returns a meaningful
diagnostic rather than wandering to \texttt{max\_scba\_iter} and
returning garbage.

\begin{table}[ht]
  \centering
  \begin{tabular}{lll}
    \hline
    option & default & legacy value \\
    \hline
    \texttt{auto\_extend\_fmax}       & \texttt{True}      & --- \\
    \texttt{fmax\_margin}             & $1.05$             & --- \\
    \texttt{zero\_mode\_projection}   & \texttt{True}      & \texttt{False} \\
    \texttt{solver}                   & \texttt{"anderson"} (derived) & --- \\
    \texttt{anderson\_safeguard}      & \texttt{True}      & \texttt{False} \\
    \texttt{gate\_on\_conservation}   & \texttt{False}     & \texttt{True} \\
    \texttt{divergence\_guard}        & \texttt{True}      & \texttt{False} \\
    \texttt{causality\_projection}    & \texttt{False}     & --- \\
    \texttt{vertex\_scale}            & $1.0$              & --- \\
    \texttt{mixing}                   & $0.3$ linear       & $0.5$ linear \\
    \texttt{eta\_factor}              & $1.0$
       ($\eta\sim\Delta\omega$) & $0.05$ \\
    \texttt{max\_scba\_iter}          & $\ge 80$           & --- \\
    \texttt{scba\_tol}                & $10^{-3}$          & --- \\
    cutoffs                           & \texttt{None}      & --- \\
    \texttt{retarded}                 & \texttt{"fft"}     & --- \\
    \hline
  \end{tabular}
  \caption{Solver options that affect convergence. The \emph{legacy
  value} column gives the setting that reproduces the previous (now
  retired) behaviour bit-for-bit; the default column is the new safe
  setting. The single most important fix is
  \texttt{auto\_extend\_fmax} (no legacy --- the previous code did
  not check the convolution-support condition).}
  \label{tab:solver_recommended}
\end{table}

% REVIEW: wrap heading math in \texorpdfstring for PDF bookmarks/intcalc
\subsection[Joint eta->0 and lambda->1 convergence diagnostic]{Joint $\eta\!\to\!0$ and $\lambda\!\to\!1$ convergence diagnostic}
\label{sec:phonon_solver:eta_lambda_diag}

The SCBA fixed-point observable, e.g.\ the anharmonic conductance
$G_{\rm anh}$, depends on two parameters that are easy to confuse with
each other:
\begin{itemize}
\item the \emph{numerical regulator} $\eta_{w}$, the imaginary part
added to the Dyson denominator
$z^{2}=(\omega+i\eta_{w})^{2}$, which controls the broadening of the
phonon spectral function and the regularisation of the retarded
Hilbert transform;
\item the \emph{physical coupling strength} $\lambda$ (\texttt{vertex\_scale}
in \texttt{transmission\_finite}), which multiplies the cubic
force-constant tensor $\Phi^{(3)}\to\lambda\Phi^{(3)}$, so that
$\Sigma\propto\lambda^{2}$
\cite{mera_inelastic_2012}.
\end{itemize}
The physical observable lives at $(\eta_{w}\!\to\!0,\lambda\!=\!1)$;
the two natural routes to it from a numerically tractable point are the
$\eta_{w}\!\to\!0$ extrapolation and the $\lambda\!\to\!1$ Pad\'e
continuation described below, and the diagnostic content lies in
whether they agree.

\paragraph{$\eta_{w}\!\to\!0$ at fixed $\lambda=1$.}
Sweeping $\eta_{w}\in\{\eta_{1},\eta_{2},\ldots\}$ from a safely large
value (large enough to broaden any anharmonically renormalised soft
mode, i.e.\ $\eta_{w}\!\gtrsim\!\sqrt{\,|\min\lambda(H_{D}+\mathrm{Re}\,\Sigma^{R})|\,}$\!\!,
see \texttt{dynamical\_stability\_diagnostic}) down to the
limit of stable SCBA convergence, and extrapolating
$G_{\rm anh}(\eta_{w})$ linearly to $\eta_{w}\!=\!0$, isolates the
regulator-free SCBA prediction at the physical coupling.

\paragraph{$\lambda\!\to\!1$ at fixed safe $\eta_{w}$.}
The lowest-order-approximation rescaling
of~\cite{mera_inelastic_2012} runs the SCBA at
$\lambda\in(0,1)$, where the contraction
$|\Sigma|/|H|\propto\lambda^{2}$ is small and convergence is
mechanical, and fits a Pad\'e in $\lambda^{2}$. Evaluating the fit at
$\lambda^{2}\!=\!1$ recovers an estimate of the SCBA answer at full
coupling without iterating an unstable map -- and, importantly, at
the same $\eta_{w}$ as the rest of the sweep, so this is also a
``regulated SCBA'' estimate of the same quantity.

\paragraph{Interpretation of the comparison.}
Both routes target the same underlying object at the same $\eta_{w}$:
the fixed-point self-energy of the bubble closure
$\Sigma=\frac{1}{2}\Phi G[\Sigma] G[\Sigma]\Phi$, which sits inside the
Baym--Kadanoff \(\Phi\)-derivable family
\cite{dahlen_LW_2006}: an
SCBA-equivalent closure is the functional derivative of a single
two-loop $\Phi[G]$, so the same fixed point is reachable from
either direction. Hence:
\begin{itemize}
\item \textbf{Agreement} ($|\Delta|/G^{\eta}_{\rm extr}\!\lesssim\!10^{-2}$
within the joint fit uncertainty): the SCBA at $\lambda\!=\!1$ is
consistent with its own analytic continuation in $\lambda^{2}$, the
$\lambda$-dependence of the fixed point is well represented by the
chosen Pad\'e family, and the bubble closure is internally
self-consistent. The headline $G_{\rm anh}$ can be quoted with
confidence to within the joint extrapolation error.
\item \textbf{Disagreement.} Two non-exclusive readings:
\begin{enumerate}
\item the assumed Pad\'e family (quadratic in $\lambda^{2}$, here)
is too restrictive to capture the actual $\lambda$-dependence of the
SCBA fixed point near $\lambda^{2}\!=\!1$ -- a curve-fitting
artefact, testable by re-fitting with a higher-order or a
saturating Pad\'e $[1/1]$ form $a+b\lambda^{2}/(1+c\lambda^{2})$;
\item the SCBA bubble closure misses contributions outside the
two-loop $\Phi[G]$ -- in particular the \emph{vertex correction}
(a ladder of bubbles inside the bubble), which contributes at
$\mathcal{O}(\lambda^{4})$ at the level of $\Sigma$
\cite{mera_inelastic_2012}, and the four-phonon scattering
channel, which appears at the same order in the cubic coupling once
the quartic FC tensor is included
\cite{han_fourphonon_2021}. In electron-phonon transport,
Migdal's argument parametrically suppresses vertex corrections by
the sound-to-Fermi velocity ratio
\cite{roy_migdal_2014}; for phonon-phonon there is no analogous
parametric suppression (the relevant velocities are all atomic
vibrational scales), and the gap $\Delta$ measures the cumulative
size of these channels at the SCBA fixed point.
\end{enumerate}
\end{itemize}
The two readings can be distinguished by running additional $\lambda$
points near $\lambda^{2}\!=\!1$ and re-fitting: if the Pad\'e
extrapolation converges as more points are added (or as the
$[1/1]$ saturating ansatz is used), the disagreement is a fit-family
artefact; if it does not, the residual gap is a physical statement
about missing higher-order channels.

\paragraph{Implementation.}
\texttt{phonon/scripts/scba\_convergence\_diagnostic.py} runs both
sweeps in one invocation, fits the two extrapolations, and writes a
\texttt{convergence\_assessment.json} reporting
$\mathcal{O}^{\eta}_{\rm extr}$, $\mathcal{O}^{\lambda}_{\rm Pad\acute{e}}$,
the relative gap
$|\Delta|/\mathcal{O}^{\eta}_{\rm extr}$, and a verdict
(``consistent'', ``borderline'', ``inconsistent'') against a
user-set tolerance. A four-panel plot
(\texttt{eta\_lambda\_diagnostic.pdf}) overlays $G_{\rm anh}(\eta_{w})$
with the linear fit, $G_{\rm anh}(\lambda^{2})$ with the Pad\'e fit,
the two extrapolated values with their fit-uncertainty bands, and
the spectral heat current $j(\omega)$ overlay across both sweeps so
the user can see at which frequencies the two routes diverge.

The script is tuned for affordable runs (looser \texttt{scba\_tol},
short $\eta$ and $\lambda$ lists) so it can serve as a routine
sanity check on every new device or temperature regime, not just as a
deep diagnostic. The Bescond pure-Pad\'e sweep already lives in
\texttt{scba\_diagnostic.py}; this new driver pairs it with the
$\eta$-extrapolation half of the diagnostic and writes the verdict in
machine-readable form.

\subsection{Zero modes and acoustic sum rules}

For the d5a SiNW the FC2 acoustic sum rule and the standard FC3 sum
rule (sum over all atoms of one leg) hold to machine precision, so the
acoustic branches sit at $\omega=0$ at $\Gamma$ and a uniform
translation carries no spurious 3-phonon coupling; the device
dynamical matrix has no imaginary modes. The $\omega=0$ sample is
excluded from every physical integral by \texttt{pos\_mask}, so the
$1/\omega$ divergence of the Bose factor cannot leak into the heat
current; its only residual effect is through the bubble convolution,
where the three \texttt{dc\_handling} options differ by well under a
per cent. Note that \texttt{enforce\_asr\_fc3\_matrices} (used by the
separable $q$-path under an opt-in flag) enforces a \emph{stronger}
per-sublattice sum than the physical sum-over-all-atoms ASR; the dense
path correctly does not apply it.

The bubble self-energy, however, is \emph{not} guaranteed to inherit a
discrete acoustic sum rule from an ASR-clean FC3, and the
H-terminated wire has an additional near-zero mode that the standard
ASR does not cover. The cell dynamical matrix
$H_{00}+H_{01}+H_{01}^{\dagger}$ at $\Gamma$ for d5a has four (not
three) eigenvalues below $10^{-4}\,\omega_{\max}^{2}$: the three
rigid-cartesian translations and the rigid rotation about the wire
axis (the ``twist''), the latter at
% REVIEW: n_B = 1/(exp(hf/kT)-1) ~= 223 for f=0.028 THz at 300 K (was ~700)
$\omega\approx 0.028$\,THz. The construction
$\Sigma^{\gtrless}\propto\Phi G^{\gtrless}G^{\gtrless}\Phi$ contains
two contracted Green's functions that do not annihilate this subspace,
and on the real device a generic $\mathrm{Re}\,\Sigma^{R}$ acquires a
$\sim 0.25$ weight across the four-dimensional zero subspace. At
% REVIEW: n_B = 1/(exp(hf/kT)-1) ~= 223 for f=0.028 THz at 300 K (was ~700)
$T=300$\,K the Bose factor $n_{B}\approx 223$ on the twist makes the
soft-mode $G^{<}$ enormous, and the bubble's IR contribution there
drives the SCBA non-contractive even when the grid covers
$2\omega_{\max}$. The \texttt{zero\_mode\_projection} option
(\Cref{ssec:solver_safeguards}) enforces a discrete ASR on $\Sigma$
itself by $\Sigma\leftarrow Q\Sigma Q$ with $Q$ the orthogonal
projector onto the complement of \emph{every} cell eigenmode below the
threshold --- translations plus twist plus any other accidentally-soft
mode. The leakage drops to machine zero and the bubble can no longer
renormalise the near-zero modes nor source an IR catastrophe through
them.

\subsection{Literature consistency}

% REVIEW: replace hard-typed citation with \cite{guo_quantum_2020}
The bubble is the Migdal form of Guo \emph{et al.}~\cite{guo_quantum_2020}, Eq.~8; the prefactor, the
$\omega$-convolution and
the Kramers--Kronig reconstruction all match the theory note. The
scheme is the self-consistent Born approximation --- iteration zero is
the lowest-order Born self-energy and the converged fixed point is the
SCBA --- with 3-phonon scattering only, consistent with reference NEGF
phonon-transport implementations. The Si/Ge interface benchmark data
% REVIEW: "Guo, Latour and Tian" attribution does not match the Guo bib entry (no Latour/Tian coauthor); use \cite
of Guo \emph{et al.}~\cite{guo_quantum_2020} ship in
\texttt{phonon/examples/literature\_fig5b.npz} for the overlay
produced by \texttt{si\_ge\_interface\_quatrex.py}.
